\documentclass[11pt]{article}
\usepackage[margin=1.1in]{geometry}
\usepackage{amsmath,amssymb,amsthm}
\usepackage{booktabs}
\usepackage[hidelinks]{hyperref}
\usepackage{microtype}
\input{numbers.tex}

\newtheorem{theorem}{Theorem}
\newtheorem{lemma}[theorem]{Lemma}
\newtheorem{corollary}[theorem]{Corollary}
\newtheorem*{remark}{Remark}
\newcommand{\vphi}{\varphi}

\title{The twelve-block colouring is optimal:\\ $\delta_3=117/2192$}
\author{Carlos Toledo\thanks{Independent researcher, \href{https://carlostoledo.co}{carlostoledo.co}; contact: carlos@carlostoledo.co.
Preprint v0.1. The method, the certificates and this text were produced by the author working with Claude (Anthropic).
Every certificate is re-checked by an independently written verifier in exact rational arithmetic. The proof is
computer-assisted and has not yet been refereed. [OWNER: affiliation and disclosure wording; decide on release.]}}
\date{October 2026}

\begin{document}
\maketitle

\begin{abstract}
Let $V(n)$ be the least number of monochromatic three-term arithmetic progressions in a
two-colouring of $\{1,\dots,n\}$, and $\delta_3=\liminf V(n)/n^2$. Graham offered \$100 for its
value (Erd\H{o}s problem \#1186). Parrilo, Robertson and Saracino (2008) proved
$1675/32768\le\delta_3\le117/2192$. Their upper bound comes from a colouring by twelve blocks,
which has been conjectured to be optimal. We prove that it is:
$\delta_3=117/2192$. The key statement is that the quadratic form $Q(\vphi)=\iint_{a/2\le b\le(1+a)/2}\vphi(a)\vphi(b)$
satisfies $Q\ge-5/137$ on all measurable $\vphi:[0,1]\to[-1,1]$, and the twelve-block function and its
negative are the only minimisers.

The proof expands $Q$ around the conjectured optimum. In these coordinates the first-order
term is an explicit nonnegative function $h$. A discretisation built on $h$ loses nothing at
first order. The space of colourings is then cut into five regions, and each region is closed by
an exact certificate. Near the optimum the certificate has zero gap. Away from it, the
certificates leave a margin of at least $\MinSlab$. All five certificates are re-verified by an
independent program in rational arithmetic.
\end{abstract}

\section{Introduction}
Colour $1,\dots,n$ red and blue. How few monochromatic progressions $\{a,a+d,a+2d\}$ ($d\ge1$)
can there be? A random colouring gives $n^2/16$. Parrilo, Robertson and Saracino \cite{PRS} found
a better colouring. Butler, Costello and Graham \cite{BCG} later showed how to find it by experiment
and found block colourings that beat random for longer progressions. The colouring of \cite{PRS}
splits $[n]$ into twelve blocks of relative lengths
\[28,\ 6,\ 28,\ 37,\ 59,\ 116,\ 116,\ 59,\ 37,\ 28,\ 6,\ 28\qquad(\text{out of }548),\]
alternating in colour, and gives $V(n)\le\frac{117}{2192}n^2(1+o(1))$. The best published lower
bound has been $1675/32768\approx0.0511$ \cite{PRS}. Greenwood, Kariv and Williams \cite{GKW} proved that the
twelve-block colouring is optimal among antisymmetric colourings with at most twelve blocks.

\begin{theorem}\label{thm:main}
$\displaystyle\delta_3=\lim_{n\to\infty}\frac{V(n)}{n^2}=\frac{117}{2192}.$
\end{theorem}

\section{The reduction}
We use the counting lemma of \cite{PRS}. Let $T$ be the number of bichromatic pairs $a<b$
with $a\equiv b\pmod 2$, and $N^+$ the number of ordered bichromatic pairs $(u,v)$ with
$2v-u\in[n]$. Then
\[2V=\tfrac{n^2}{2}-|N^+|-|T|+O(n),\qquad |T|\le \tfrac{n^2}{8}+O(n).\]
A non-monochromatic progression has exactly two bichromatic pairs. The outer pair is counted
in $T$, and every inner pair is counted in $N^+$. For the bound on $T$, write
$T=r_ob_o+r_eb_e$ by parity classes.

For a colouring, let $\vphi:[0,1]\to\{\pm1\}$ be its step function. Then
$|N^+|=n^2\iint_R\frac12(1-\vphi(a)\vphi(b))\,da\,db+O(n)$, where
\[R=\{(a,b)\in[0,1]^2:\ a/2\le b\le (1+a)/2\},\qquad |R|=\tfrac12 .\]
Put $Q(\vphi)=\iint_R\vphi(a)\vphi(b)$. Then
\begin{equation}\label{eq:red}
\frac{V(n)}{n^2}\ \ge\ \frac1{16}+\frac{Q^*}{4}-O\!\left(\frac1n\right),\qquad
Q^*=\inf\{Q(\vphi):\ \vphi:[0,1]\to[-1,1]\ \text{measurable}\}.
\end{equation}
Let $\vphi^*$ be the twelve-block function, with $\vphi^*=+1$ on the first block. Then
$Q(\vphi^*)=-5/137$, and $\frac1{16}-\frac5{548}=\frac{117}{2192}$. So
Theorem~\ref{thm:main} follows from \cite{PRS} together with:

\begin{theorem}\label{thm:A}
$Q(\vphi)\ge-5/137$ for every measurable $\vphi:[0,1]\to[-1,1]$. Equality holds only for
$\vphi=\pm\vphi^*$ almost everywhere.
\end{theorem}

\section{Expanding around the twelve blocks}
Let $K$ be the symmetrised kernel of $R$, so that $Q(\vphi)=\langle\vphi,K\vphi\rangle$. Put
$g^*=K\vphi^*$:
\[g^*(a)=\tfrac12\Bigl(\int_{a/2}^{(1+a)/2}\vphi^*+\int_{[2a-1,2a]\cap[0,1]}\vphi^*\Bigr),\qquad h=-\vphi^*g^*.\]
Any $\vphi$ with values in $[-1,1]$ can be written as $\vphi=\vphi^*(1-2\sigma)$ with
$\sigma=(1-\vphi\vphi^*)/2\in[0,1]$. Here $\sigma$ is the part of $\vphi$ that is flipped against
$\vphi^*$, and $\int\sigma=\frac12\|\vphi-\vphi^*\|_1$. Expanding the quadratic form gives the exact
identity
\begin{equation}\label{eq:E}
E(\sigma):=\frac{Q(\vphi)-Q(\vphi^*)}{4}=\int_0^1h\,\sigma+Q(\vphi^*\sigma).
\end{equation}

\begin{lemma}\label{lem:h}
The function $h$ is continuous and piecewise linear, with nodes at the multiples of $1/1096$.
It satisfies $h\ge0$, and it vanishes exactly at the eleven breakpoints of $\vphi^*$, where its
one-sided slopes are $1$ or $3/2$. Moreover $\int_0^1h=5/137$.
\end{lemma}
\begin{proof}
The kinks of $g^*$ occur only where $a/2$, $(1+a)/2$, $2a$ or $2a-1$ meets an edge of $\vphi^*$,
or at $a=\frac12$. Every edge is a multiple of $2/1096$, so every kink is a multiple of $1/1096$.
The values at the nodes are computed exactly. At each edge $g^*=0$, so $h$ is continuous there.
\end{proof}

Since $\vphi\mapsto-\vphi$ sends $\sigma$ to $1-\sigma$ and does not change $Q$, it suffices to prove
$E(\sigma)\ge0$ when $m:=\int\sigma\le\frac12$.

\section{Discretisation without first-order loss}
Partition $[0,1]$ into cells $I_i$ whose endpoints include every edge of $\vphi^*$. On $I_i$,
$\vphi^*$ is a constant $s_i$; write $w_i=|I_i|$ and $y_i=\frac1{w_i}\int_{I_i}\sigma\in[0,1]$.

\begin{lemma}[bathtub]\label{lem:bath}
Let $a_i=w_i\min_{I_i}h$, and let $m_i(v)=|\{a\in I_i: h(a)<v\}|$. Put
$b_i=\frac{w_i^2}{2}\big/\sup m_i'$, or $b_i=0$ if $h$ is constant on some piece of $I_i$. Then
\[\int_{I_i}h\sigma\ \ge\ a_iy_i+b_iy_i^2\qquad\text{whenever }0\le\sigma\le1.\]
\end{lemma}
\begin{proof}
For a fixed average $y$, the integral is smallest when $\sigma$ fills the lowest levels of $h$.
So it is at least $\beta(y)=\int_0^{w_iy}h^{\uparrow}$, where $h^{\uparrow}$ is the increasing
rearrangement of $h$ on $I_i$. We have $\beta(0)=0$ and $\beta'(0)=a_i$. Also
$\beta''(y)=w_i^2/m_i'(h^\uparrow(w_iy))\ge 2b_i$, and the claim follows.
\end{proof}

\begin{lemma}\label{lem:disc}
$E(\sigma)\ge L(y)$, where
\begin{align*}
L(y)={}&\sum_i\bigl(a_iy_i+b_iy_i^2\bigr)+\sum_{I_i\times I_j\subset R}s_is_jw_iw_jy_iy_j\\
&+\sum_{\text{cut},\,s_is_j=1}\max\bigl(0,\ w_iw_jy_iy_j-A^{\rm out}_{ij}\bigr)
+\sum_{\text{cut},\,s_is_j=-1}\max\bigl(-A^{\rm in}_{ij},\ -w_iw_jy_iy_j\bigr).
\end{align*}
Here a cut pair is one where $\partial R$ crosses $I_i\times I_j$, and
$A^{\rm in},A^{\rm out}$ are the areas inside and outside $R$.
\end{lemma}
\begin{proof}
On a pair inside $R$, the integral of $\sigma(a)\sigma(b)$ factorises. On a cut pair,
$0\le\iint_{\rm in}\sigma\sigma\le\min(A^{\rm in},w_iw_jy_iy_j)$, because $0\le\sigma\le1$.
Also $\iint_{\rm in}\sigma\sigma=\iint_{\rm full}\sigma\sigma-\iint_{\rm out}\sigma\sigma\ge w_iw_jy_iy_j-A^{\rm out}$.
\end{proof}

$L(0)=0$, and the first-order part of $L$ at $y=0$ is exact up to the bathtub. This is the
difference from the usual discretisation in $\vphi$ \cite{PRS}. There, every cut pair loses a
first-order amount, which caps any certified bound strictly below $117/2192$.

\section{Five certificates}
Write $m=\sum_iw_iy_i=\int\sigma$. We cover $m\in[0,\frac12]$ by five regions. Each region has
its own grid and an exact certificate (Table~\ref{tab}). The two grids are:
\begin{itemize}
\item 130 cells: multiples of $24/1096$, steps of $4/1096$ within $8/1096$ of every edge, and
their mirror images;
\item 146 cells: multiples of $8/1096$ together with the edges.
\end{itemize}

\paragraph{Slabs $\mathrm{lo}\le m\le \mathrm{hi}$.}
We use an exact polynomial identity in $v=(1,y)$:
\[L_\theta(y)-B=v^\top Sv+\sum_f\mu_ff(y),\qquad \mu\ge0,\ \theta\in[0,1],\ S+\varepsilon I\succeq0 .\]
\begin{itemize}
\item $L_\theta\le L$ replaces each maximum in Lemma~\ref{lem:disc} by a convex combination.
\item Each form $f$ is one of two kinds. Either it is a product of two of the factors $y_i$,
$1-y_i$, $m-\mathrm{lo}$, $\mathrm{hi}-m$, each nonnegative on the region; or it is a
triangle inequality $1+s_1z_iz_j+s_2z_jz_k+s_3z_iz_k$ with $s_1s_2s_3=1$ and $z=1-2y$.
\end{itemize}
Hence $L\ge B-\varepsilon(1+n)$ on the region. The products with $m-\mathrm{lo}$ and $\mathrm{hi}-m$
matter. They stop the relaxation from averaging the two optima $\pm\vphi^*$, and from averaging
points of different regions.

\paragraph{The region $m\le\Rin$.}
Here every term vanishes at $y=0$:
\begin{multline*}
L_0(y)=y^\top Py+\sum_{i,j}\lambda_{ij}\,y_i(1-y_j)+\sum_j\nu_j\,y_j(\Rin-m)\\
+\kappa\,m(\Rin-m)+\sum_{i,j}N_{ij}\,y_iy_j+\sum_ic_iy_i ,
\end{multline*}
with $P\succ0$ and all multipliers nonnegative. $L_0\le L$ uses the cut-pair bounds $0$ and
$-w_iw_jy_iy_j$. This certificate is sharp: the forms $y_j(1-y_i)$ convert the linear budget
$a_iy_i$ into exactly the bathtub's quadratic cost, because
$\sum_j jy_j-\sum_{i<j}y_iy_j=\sum_{i<j}y_j(1-y_i)$. The breakpoint Hessian of $Q$ at $\vphi^*$
is positive definite, with smallest eigenvalue $0.1266$, and this supplies the margin.

\begin{table}[h]\centering
\begin{tabular}{cccccc}\toprule
region $m=\int\sigma$ & cells & bound & nonzero multipliers & triangles\\\midrule
\CoverTable
\bottomrule\end{tabular}
\caption{The cover of $m\in[0,\frac12]$. The bounds are the exact certified values, truncated.}\label{tab}
\end{table}

\begin{proof}[Proof of Theorem~\ref{thm:A}]
The five regions cover $m\in[0,\frac12]$, so $E\ge0$ everywhere. If $y\neq0$ and $m\le\Rin$, then
$E\ge L_0(y)\ge y^\top Py>0$. If $y=0$, then $\sigma=0$ almost everywhere. Every slab has $E>0$.
So $E=0$ only at $\sigma=0$, that is $\vphi=\vphi^*$. The case $m\ge\frac12$ follows by symmetry.
\end{proof}

\begin{corollary}
$\vphi^*$ minimises $Q$ on the whole ball $\|\vphi-\vphi^*\|_1\le2\cdot\Rin$, fractional
perturbations included. Moreover, every $\vphi$ with $\frac12\|\vphi-\vphi^*\|_1\ge\Rin$ and
$\frac12\|\vphi+\vphi^*\|_1\ge\Rin$ has $Q(\vphi)\ge-\frac5{137}+4\cdot\MinSlab$.
\end{corollary}

\section{Verification}
The multipliers come from floating-point semidefinite programs (CVXPY with Clarabel).
They are rounded to rationals. The matrix $S$ or $P$ is then recomputed exactly and
positive definiteness is proved exactly. A second program, written separately, re-checks
everything. It:
\begin{itemize}
\item re-derives $\vphi^*$, $h$ and its linearity;
\item recomputes every bathtub minorant, and every cell-pair area by polygon clipping;
\item rebuilds every form from its name;
\item checks the identity~\eqref{eq:E} and Lemma~\ref{lem:disc} in exact arithmetic on random step functions;
\item proves definiteness by fraction-free elimination (Sylvester's criterion);
\item checks that the regions cover $[0,\frac12]$.
\end{itemize}
Forged certificates are rejected, each on its own line: a bound raised by $10^{-3}$, the triangle
multipliers dropped, the cut combinations altered, a grid missing a breakpoint, a widened slab, a
doubled inner radius, an overspent linear budget, the inner RLT multipliers dropped, and a missing
slab. One-line edits of the verifier itself are caught as well, by the genuine set or by planted
forgeries. Nine of ten such edits are caught; the tenth is provably harmless.
The certificate files, both programs and their SHA-256 hashes accompany this preprint.
[OWNER: repository / DOI.]

\section{Remarks}
The proof does not use the parity structure: \eqref{eq:red} only needs $|T|\le n^2/8$.
The argument is computer-assisted. Its weight rests on Lemmas~\ref{lem:h}--\ref{lem:disc},
which are elementary, and on five finite certificates, which anyone can re-check in about
forty seconds. We would welcome an independent re-verification.

\begin{thebibliography}{9}
\bibitem{PRS} P.~A.~Parrilo, A.~Robertson, D.~Saracino, On the asymptotic minimum number of
monochromatic 3-term arithmetic progressions, \emph{J. Combin. Theory Ser. A} 115 (2008); arXiv:math/0609532.
\bibitem{BCG} S.~Butler, K.~P.~Costello, R.~Graham, Finding patterns avoiding many monochromatic
constellations, \emph{Experiment. Math.} 19 (2010) 399--411.
\bibitem{GKW} T.~Greenwood, J.~Kariv, N.~Williams, From discrete to continuous: monochromatic
3-term arithmetic progressions, arXiv:2301.00336.
\bibitem{EP} T.~F.~Bloom, Erd\H{o}s problem \#1186, \url{https://www.erdosproblems.com/1186}.
\end{thebibliography}
\end{document}
